\section{Dafny: Technical Audit}\label{sec:audit}

\footnote{also look at broken / incorrect examples and see equivalence, they should be}

\footnote{highlighting important parts of examples}
\footnote{any limitations that only arise with more complex data structures?}

There are many cases where the user needs to provide assistance to Dafny in order to prove equivalence. In this section, I will provide a series of examples to demonstrate where Dafny's limitations are in this area, and how the user can provide assistance to allow Dafny to overcome these limitations and prove equivalence. These limitations include a series of problems when generating and proving auxiliary lemmas, problems with proving termination of the equivalence lemmas in the presense of termination proofs for the individual functions, and problems with proving the equivalence of Dafny methods with imperative features, such as loops.

The majority of the following examples have been taken from the program equivalence tests in the GitHub repository for the Stainless Scala verification tool \parencite{StainlessFrontendsBenchmarks}. I converted all of their tests into Dafny functions, and proved their equivalence. Many of these tests consisted of more than 2 functions, each in their own file. For these cases, I concatenated all of these functions into a single Dafny file to form a single test, and proved the equivalence of each function with a single function, thus proving, through the transitivity of equivalence, that all of these functions were pairwise equivalent. This resulted in a total of 62 tests.

These tests cover a great variety of types of programs. The number of lines of code in each test ranges from around 10 to over 200. They cover a variety of inbuilt and custom data structures, including sequences, lists and the Option datatype. Some tests also cover concepts such as mutual recursion, generic types and higher-order functions.

Using these tests, I have attempted to characterise the cases where Dafny required user assistance. In this way, I have identified 7 different types of limitations. \Cref{tbl:limitations-summary} summarises the number of tests that demonstrate each type of limitation. Many examples demonstrate multiple types of limitations, therefore, these are counted multiple times. This table also includes imperative equivalence tests that I wrote myself. All of the tests from the Stainless repository can be proven equivalent by Stainless fully automatically, except for the single test involving higher-order functions. 

\begin{table}[ht]
  \centering
  \begin{tabular}{ |c|c|c| } 
  \hline
  Type of Limitation & Number of Examples & Source \\
  \hline
  None (Fully Automatic) & 19 & \parencite{StainlessFrontendsBenchmarks} \\
  \hline
  \nameref{sec:counterexamples} & 1 & \parencite{StainlessFrontendsBenchmarks} \\ 
  \hline
  \nameref{sec:equivalence-lemmas} & 17 & \parencite{StainlessFrontendsBenchmarks} \\
  \hline 
  \nameref{sec:auxiliary-lemmas} & 3 & \parencite{StainlessFrontendsBenchmarks} \\
  \hline 
  \nameref{sec:induction} & 4 & \parencite{StainlessFrontendsBenchmarks} \\
  \hline
  \nameref{sec:termination} & 17 & \parencite{StainlessFrontendsBenchmarks} \\
  \hline
  \nameref{sec:imperative-equivalence} & 7 & New \\
  \hline
  \nameref{sec:higher-order} & 1 & \parencite{StainlessFrontendsBenchmarks} \\
  \hline
  \end{tabular}
  \caption{Summary of the number of examples that demonstrate each type of limitation}
  \label{tbl:limitations-summary}
\end{table}

The following sections describe each type of limitation that I have identified. Each section contains at least one example to clearly demonstrate the limitation being discussed. I have tried to pick the simplest, minimal examples that I have access to, to remove any unnecessary difficulty. 

\subsection{Necessary Datastructures}

I will require the following definition of a list data structure:
\begin{lstlisting}[caption=List Data Structure]
datatype List<T> = Nil | Cons(head: T, tail: List<T>)
\end{lstlisting}

\subsection{Incorrect Counterexamples}\label{sec:counterexamples} 
In this example, I will provide two functions that attempt to determine whether a list of integers is sorted in non-decreasing order. The first function is as follows, involving the main function \lstinline|isSortedR| and a helper function \lstinline|loop|:

\begin{lstlisting}[caption=\lstinline|isSortedR| and \lstinline|loop| Functions, label=lst:isSortedR]
function loop(p: int, l: List<int>): bool
  decreases l
{
  match l
  case Nil => true
  case Cons(x, xs) => p <= x && loop(x, xs)
}

function isSortedR(l: List<int>): bool
{
  match l
  case Nil => true
  case Cons(x, xs) => loop(x, xs)
}
\end{lstlisting}
%
The second function is as follows, involving the main function \lstinline|isSortedA| and helper functions \lstinline|iter| and \lstinline|leq|:

\begin{lstlisting}[caption=\lstinline|isSortedA| function and \lstinline|iter| and \lstinline|leq| helper functions]
function leq(cur: int, next: int): bool
{
  cur < next
}

function iter(l: List<int>): bool
{
  match l
  case Nil => true
  case Cons(x, Nil) => true
  case Cons(x, Cons(y, ys)) => leq(x, y) && iter(Cons(y, ys))
}

function isSortedA(l: List<int>): bool
{
  match l
  case Nil => true
  case Cons(_, Nil) => true
  case Cons(x, Cons(y, ys)) => x <= y && iter(Cons(y, ys))
}
\end{lstlisting}
%
These functions are not equivalent, as \lstinline|isSortedA| uses a strictly less-than comparison in the \lstinline|leq| function, when it should use a less-than-or-equal-to comparison. Therefore, Dafny is unable to prove the following equivalence lemma:
\begin{lstlisting}[caption=Equivalence Lemma for \lstinline|isSortedA| and \lstinline|isSortedR|]
lemma isSortedEquivalence(l: List<int>)
  ensures isSortedA(l) == isSortedR(l)
{}
\end{lstlisting}
%

\footnote{Why does it provide an incorrect counter example?}

However, it incorrectly provides \lstinline|Cons(-38, Cons(7681, Nil))| as a counterexample. I extended the above proof to include and use a helper lemma, as shown below:
\begin{lstlisting}[caption=Equivalence Lemma with Helper for \lstinline|isSortedA| and \lstinline|isSortedR|]
lemma isSortedEquivalenceHelper(x: int, xs: List<int>)
  decreases xs
  ensures iter(Cons(x, xs)) == loop(x, xs)
{}

lemma isSortedEquivalence(l: List<int>)
  ensures isSortedA(l) == isSortedR(l)
{
  match l
  case Nil => {}
  case Cons(x, Nil) => {}
  case Cons(x, Cons(y, ys)) => isSortedEquivalenceHelper(y, ys);
}
\end{lstlisting}
%

\footnote{Possible extension: fixing counter examples? is this something I can do something about?}

Now, Dafny is able to provide a correct counterexample of \lstinline|x = 8855| and \lstinline|xs = Cons(8855, Nil)| for the helper lemma. However, though this is the simplest counterexample in terms of list length, it is not the simplest example in terms of the integers used, as \lstinline|x = 0| and \lstinline|xs = Cons(0, Nil)| would be sufficient.

As asserted by \textcite{putrueleMaskDToolMeasuring2022}, since there are multiple translations steps between the Dafny source code and the verification conditions sent to the Z3 SMT solver, the counterexamples produced by it can be difficult to read and understand. They also note these counterexamples may sometimes be incorrect, due to the Z3 solver being incomplete.

\subsection{Equivalence Lemmas for Helper functions}\label{sec:equivalence-lemmas}

For this example, I will provide two solutions to the problem of reducing a list down its unique elements. Both solutions each require two helper functions, where each of the three pairs of functions across the two solutions are equivalent. The first pair of functions are \lstinline|isin| and \lstinline|find|. These functions determine whether a given element is in a list:
\begin{lstlisting}[caption=\lstinline|isin| and \lstinline|find| Functions]
function find(lst: List<int>, n: int): bool
{
  match lst
  case Nil => false
  case Cons(hd, tl) => n == hd || find(tl, n)
}

function isin(a: int, lst: List<int>): bool
{
  match lst
  case Nil => false
  case Cons(hd, tl) => a == hd || isin(a, tl)
}
\end{lstlisting}
%
Dafny requires that these helper functions are explicitly proven equivalent to be able to prove that the main functions are equivalent. Therefore, the user must write the following equivalence lemma:
\begin{lstlisting}[caption=Equivalence Lemma for \lstinline|isin| and \lstinline|find|]
lemma findEquivalence(l: List<int>, n: int)
  ensures find(l, n) == isin(n, l)
{}
\end{lstlisting}
%
\textcite{milovancevicProvingDisprovingEquivalence2023} were able to automatically generate these helper lemmas by preprocessing the two solutions to find helper functions with compatible signatures: where the arguments of one function can be permuted so that their types match the types of the arguments of another function. Once candidate pairs of functions are found, equivalence lemmas such as the one above can be automatically generated and proven. If multiple functions share the same signature, test cases are used to narrow down the search space. These can be defined by the user or generated from counterexamples found when disproving equivalence of other incorrect candidate submissions.

The second pair of functions are \lstinline|unique| and \lstinline|distinct|. These functions take two lists as input, and compute the unique elements of their union:

\begin{lstlisting}[caption=\lstinline|unique| and \lstinline|distinct| Functions]
function unique(l: List<int>, r: List<int>): List<int>
{
  match l
  case Nil => r
  case Cons(hd, tl) =>
    if !find(r, hd) then unique(tl, append(r, Cons(hd, Nil)))
    else unique(tl, r)
}

function distinct(a: List<int>, b: List<int>): List<int>
  decreases b
{
  match b
  case Nil => a
  case Cons(hd, tl) =>
    if isin(hd, a) then distinct(a, tl)
    else distinct(append(a, Cons(hd, Nil)), tl)
}
\end{lstlisting}
%

To prove their equivalence, the user must write the following equivalence lemma:
\begin{lstlisting}[caption=Equivalence Lemma for \lstinline|unique| and \lstinline|distinct|]
lemma uniqEquivalenceHelper(l: List<int>, r: List<int>)
  ensures unique(l, r) == distinct(r, l)
{
  match l
  case Nil => {}
  case Cons(hd, tl) => findEquivalence(r, hd);
}
\end{lstlisting}
%

Therefore, since the proof of the equivalence of \lstinline|unique| and \lstinline|distinct| relies on the equivalence of \lstinline|find| and \lstinline|isin|, the user must explicitly call the equivalence lemma for the helper functions. This lemma must be called with the correct arguments and under the correct conditions, within the main equivalence lemma (as demonstrated above).

In addition, this lemma cannot be so easily generated using the method described above, as the arguments of \lstinline|unique| and \lstinline|distinct| have the same types, but are used in different ways. Therefore, to successfully generate this lemma, the functions themselves must be processed to identify how each argument is used.

Finally, the third pair of functions \lstinline|uniqR| and \lstinline|uniqA| to find the unique elements of a list are as follows:

\begin{lstlisting}[caption=\lstinline|uniqR| and \lstinline|uniqA| Functions]
function uniqR(lst: List<int>): List<int>
{
  unique(lst, Nil)
}

function uniqA(lst: List<int>): List<int>
{
  distinct(Nil, lst)
}
\end{lstlisting}
%
These functions can be proven equivalent using the following equivalence lemma:
\begin{lstlisting}[caption=Equivalence Lemma for \lstinline|uniqR| and \lstinline|uniqA|]
lemma uniqEquivalence(l: List<int>)
  ensures uniqR(l) == uniqA(l)
{
  uniqEquivalenceHelper(l, Nil);
}
\end{lstlisting}
%

Therefore, to prove the equivalence of two functions that both use helper functions, the user must first prove the equivalence of the helper functions, and then explicitly call the corresponding lemmas within the main equivalence lemma, ensuring that all required conditions are met at the call site.

\subsection{Generating Auxiliary Lemmas}\label{sec:auxiliary-lemmas}

Sometimes, to prove the equivalence of two functions, Dafny requires helper lemmas that will be difficult to generate automatically, since they require an understanding of what the functions being compared do. For example, consider the following two functions \lstinline|limit3_1| and \lstinline|limit3_2|, which attempt to compute the triangular number of a given integer \lstinline|n| (the functions return \lstinline|n| when n is negative):
\begin{lstlisting}[caption=\lstinline|limit3_1| and \lstinline|limit3_2| Functions]
function limit3_1(n: int): int {
  if (n <= 1) then n
  else n + limit3_1(n-1)
}

function limit3_2(n: int): int {
  if (n <= 1) then n
  else
    var r := limit3_2(n-1);
    if (r >= 0) then n + r
    else r
}
\end{lstlisting}
%
The main difference between these two functions is in the \lstinline|else| branch of the \lstinline|limit3_2| function, where it will either output \lstinline|n + r| or \lstinline|r|, depending on whether \lstinline|r >= 0|. However, to reach this \lstinline|else| branch, we know that \lstinline|n| must be greater than or equal to 2, so \lstinline|n - 1| must be at least 1. Therefore, \lstinline|limit3_2(n-1)| will output a triangular number, which must be non-negative. This means that we should be able to prove that \lstinline|r| is always non-negative, so \lstinline|n + r| will always be returned. Dafny requires some help to realise this, and to prove equivalence:

\begin{lstlisting}[caption=Auxiliary and Equivalence Lemmas for the \lstinline|limit3_1| and \lstinline|limit3_2| Functions]
lemma equivalenceLimit3_Helper(n: int)
  requires n >= 1
  ensures limit3_2(n) > 0
{}

lemma equivalenceLimit3(n: int)
  ensures (limit3_1(n) == limit3_2(n))
{
  if n >= 2 { equivalenceLimit3_Helper(n); }
}
\end{lstlisting}
%
Therefore, Dafny is able to prove the equivalence of these two functions, but only after the user provides a helper lemma that would be difficult for Dafny to automatically generate. Here is another, more complicated example. Consider the following two functions \lstinline|m1| and \lstinline|m2|, which output the fibonnaci number of a given integer \lstinline|n|. However, both functions contain an extra boolean argument: 
\begin{lstlisting}[caption=\lstinline|m1| and \lstinline|m2| Functions]
function m1(n: int, flag: bool): int {
  if (n < 1) then 0
  else if (n == 1) then 1
  else m1(n - 1, !flag) + m1(n - 2, !flag)
}

function m2(n: int, mode: bool): int {
  if (n < 1) then 0
  else if (n == 1 || n == 2) then 1
  else
    var results := if (mode) then
                      m2(n-2, !mode) + m2(n-2, !mode) + m2(n-3, !mode) 
                   else
                      m2(n-1, !mode) + m2(n-2, !mode);
    results
}
\end{lstlisting}
%
The first function does not use the boolean argument at all/ The second function uses it in the \lstinline|else| branch to either compute the standard fibonacci expression \lstinline|m2(n-1, !mode) + m2(n-2, !mode)|, or the expression with the \lstinline|m2(n-1, !mode)| call unrolled once (\lstinline|m2(n-2, !mode) + m2(n-2, !mode) + m2(n-3, !mode)|). Dafny only requires assistance in proving that the unrolled function call in \lstinline|m2| is equivalent to the standard fibonacci expression. Therefore, the user must provide the following auxiliary lemma:

\begin{lstlisting}[caption=Auxiliary Lemma for \lstinline|m2|]
lemma m2_equivalence_fib(n: int)
  requires n >= 2
  ensures m2(n, true) == m2(n - 1, false) + m2(n - 2, false)
{}
\end{lstlisting}
%
Using this, Dafny can now prove the equivalence of \lstinline|m1| and \lstinline|m2|:
\begin{lstlisting}[caption=Equivalence Lemma for \lstinline|m1| and \lstinline|m2|]
lemma equivalence(n: int, flag: bool)
  ensures m1(n, flag) == m2(n, flag)
{
  if n >= 2 { m2_equivalence_fib(n); }
}
\end{lstlisting}
%
Thus, this is another example of where Dafny can prove the equivalence of two functions, but only after the user provides helper lemmas that would be difficult for Dafny to automatically generate.

\subsection{Explicitly enforcing induction to prove lemmas}\label{sec:induction}

Dafny often proves lemmas through induction. However, sometimes, it requires the user to explicitly state how to perform this induction by calling the lemma the user is trying to prove recursively with `smaller' arguments. For example, consider the two functions \lstinline|binSumM| and \lstinline|binSum1|. These functions compute the sum of two binary numbers and a carry bit, where the bits are represented as booleans, and the binary numbers are represented as lists of bits, in order of least significant to most significant. The \lstinline|binSum1| function computes the sum by listing every possible form of input, while the \lstinline|binSumM| function simplifies this by using the carry bit in a more complicated way. The functions are as follows: 

\begin{lstlisting}[caption=\lstinline|binSumM| and \lstinline|binSum1| Functions]
function binSum1(l1: List<bool>, l2: List<bool>, c: bool): List<bool>
  decreases length(l1) + length(l2)
{
  match (l1, l2, c)
    case (Nil, Nil, _) => Nil
    case (Cons(true, t1), Cons(false, t2), false) => Cons(true, binSum1(t1, t2, false))
    case (Cons(true, t1), Cons(false, t2), true) => Cons(false, binSum1(t1, t2, true))
    case (Cons(false, t1), Cons(true, t2), false) => Cons(true, binSum1(t1, t2, false))
    case (Cons(false, t1), Cons(true, t2), true) => Cons(false, binSum1(t1, t2, true))
    case (Cons(false, t1), Cons(false, t2), false) => Cons(false, binSum1(t1, t2, false))
    case (Cons(false, t1), Cons(false, t2), true) => Cons(true, binSum1(t1, t2, false))
    case (Cons(true, t1), Cons(true, t2), false) => Cons(false, binSum1(t1, t2, true))
    case (Cons(true, t1), Cons(true, t2), true) => Cons(true, binSum1(t1, t2, true))
    case (Cons(true, t1), Nil, true) => Cons(false, binSum1(t1, Nil, true))
    case (Cons(true, t1), Nil, false) => Cons(true, binSum1(t1, Nil, false))
    case (Cons(false, t1), Nil, true) => Cons(true, binSum1 (t1, Nil, false))
    case (Cons(false, t1), Nil, false) => Cons(false, binSum1(t1, Nil, false))
    case (Nil, Cons(true, t1), true) => Cons(false, binSum1(t1, Nil, true))
    case (Nil, Cons(true, t1), false) => Cons(true, binSum1(t1, Nil, false))
    case (Nil, Cons(false, t1), true) => Cons(true, binSum1(t1, Nil, false))
    case (Nil, Cons(false, t1), false) => Cons(false, binSum1(t1, Nil, false))
}

function binSumM(l1: List<bool>, l2: List<bool>, c: bool): List<bool>
  decreases length(l1) + length(l2)
{
  match (l1, l2)
    case (Nil, Nil) => Nil
    case (Cons(true, t1), Cons(false, t2)) => Cons(!c, binSumM(t1, t2, c))
    case (Cons(false, t1), Cons(true, t2)) => Cons(!c, binSumM(t1, t2, c))
    case (Cons(false, t1), Cons(false, t2)) => Cons(c, binSumM(t1, t2, false))
    case (Cons(true, t1), Cons(true, t2)) => Cons(c, binSumM(t1, t2, true))
    case (Cons(true, t1), Nil) => Cons(!c, binSumM(t1, Nil, c))
    case (Cons(false, t1), Nil) => Cons(c, binSumM(t1, Nil, false))
    case (Nil, Cons(true, t2)) => Cons(!c, binSumM(t2, Nil, c))
    case (Nil, Cons(false, t2)) => Cons(c, binSumM(t2, Nil, false))
}
\end{lstlisting}
%

Dafny can prove the equivalence of these two functions, but only if the user explicitly calls the same lemma recursively:
\begin{lstlisting}[caption=Equivalence Lemma for \lstinline|binSumM| and \lstinline|binSum1|]
lemma equivalenceBinSum(l1: List<bool>, l2: List<bool>, c: bool)
  ensures binSum1(l1, l2, c) == binSumM(l1, l2, c)
  decreases length(l1) + length(l2)
{
  match (l1, l2)
    case (Cons(_, t1), Cons(_, t2)) => equivalenceBinSum(t1, t2, c);
    case _ => 
}
\end{lstlisting}

\subsection{Providing Variants to prove termination}\label{sec:termination}

Here, I provide two implementations of a function to find the maximum element in a sequence of integers. These are the same functions as listed in \Cref{sec:dafny}. The first is as follows:
\begin{lstlisting}[caption=\lstinline|maxR| Function]
function maxR(l: seq<int>): int
{
  if |l| == 0 then -1
  else if |l| == 1 then l[0]
  else
    var m := maxR(l[1..]);
    if l[0] > m then l[0] else m
}
\end{lstlisting}
%
The second is as follows:
\begin{lstlisting}[caption=\lstinline|maxC| Function]
function maxC(l: seq<int>): int
  decreases |l|
{
  if |l| == 0 then -1
  else if |l| == 1 then l[0]
  else
    if l[0] > l[1] then
      maxC([l[0]] + l[2..])
    else
      maxC(l[1..])
}
\end{lstlisting}
%
As described in \Cref{sec:dafny}, these two functions process the list in opposite directions, but Dafny is able to prove equivalence with almost no help from the user:
\begin{lstlisting}[caption=Equivalence Lemma for \lstinline|maxC| and \lstinline|maxR|]
lemma maxEquivalence(l: seq<int>)
  ensures maxC(l) == maxR(l)
  decreases |l|
{}
\end{lstlisting}
%
Dafny requires a variant to prove termination, as it proves the equivalence lemma through induction. However, in this case, it is unable to infer this variant by itself so the user must provide it instead. This problem can be solved easily in some cases, including this one, by copying the variant from one of the functions being compared.

This solution is especially apparent if both functions have the same variant. Here I provide two implementations of a function to subtract one integer from another. The two functions are as follows:
\begin{lstlisting}[caption=\lstinline|sub| and \lstinline|mySub| Functions]
function sub(x: int, y: int): int
  decreases(if (x <= 0) then -x else x) {
  if (x == 0) then -y
  else if (x > 0) then sub(x - 1, y - 1)
  else sub(x + 1, y + 1)
}

// Computes y - x instead of x - y
function mySub(x: int, y: int): int
  decreases(if (y <= 0) then -y else y) {
  if (y == 0) then -x
  else if (y > 0) then mySub(x - 1, y - 1)
  else mySub(x + 1, y + 1)
}
\end{lstlisting}
%
The variants of these two functions are equivalent when the arguments are matched correctly. Therefore, Dafny is able to prove their equivalence as follows, but only if the user explicitly provides a variant with the same structure as those from the two functions, with the correct arguments:

\begin{lstlisting}[caption=Equivalence Lemma for \lstinline|sub| and \lstinline|mySub|]
lemma equivalenceSub(x: int, y: int)
  decreases(if (x <= 0) then -x else x)
  ensures (sub(x, y) == mySub(y, x))
{}
\end{lstlisting}
%
In addition, if at least one of the functions requires a lemma or an assertion to prove termination, this must be copied into the equivalence lemma as well. For example, here are two implementations of a function to reduce a list down to its unique elements:
\begin{lstlisting}[caption=\lstinline|solution_2| and \lstinline|uniq2| Functions]
function solution_2(lst: List<int>): List<int>
  decreases(length(lst)) {
  match lst {
    case Nil        => Nil
    case Cons(hd, tl) =>
      lemM(hd, tl);
      Cons(hd, solution_2(drop_2(tl, hd)))
  }
}

function uniq2(lst: List<int>): List<int>
  decreases(length(lst)) {
  match lst {
    case Nil => Nil
    case Cons(hd, tl) =>
      lem2(hd, tl);
      Cons(hd, uniq2(drop(hd, tl)))
  }
}
\end{lstlisting}
%
The functions \lstinline|drop_2| and \lstinline|drop| are helper functions that remove all occurrences of an integer from a list of integers. The definitions of these functions are irrelevant for this example, so I will omit them. All that is required is that they are equivalent, when the arguments are matched correctly:
\begin{lstlisting}[caption=Equivalence Lemma for \lstinline|drop_2| and \lstinline|drop|]
lemma equivalenceDrop(a: int, lst: List<int>)
  ensures (drop(a, lst) == drop_2(lst, a))
{}
\end{lstlisting}
%
To prove termination, Dafny requires the lists inputted into the recursive calls to \lstinline|solution_2| and \lstinline|uniq2| to be `smaller' than the lists provided in the original calls. Therefore, the outputs of \lstinline|drop_2| and \lstinline|drop| must be `smaller' than the original lists. Dafny cannot infer this by itself, therefore, the user must explicitly provide lemmas. \lstinline|solution_2| and \lstinline|uniq2| use the lemmas \lstinline|lemM| and \lstinline|lem2| respectively for this purpose. These lemmas are equivalent, so I will only provide \lstinline|lem2|. For this example, it is sufficient to prove one list is `smaller' than another by using their lengths:
\begin{lstlisting}[caption=\lstinline|lem2| Helper Lemma]
lemma lem2(a: int, lst: List<int>)
  ensures length(drop(a, lst)) <= length(lst)
{}
\end{lstlisting}
%
Using \lstinline|equivalenceDrop|, Dafny is able to prove the equivalence of \lstinline|solution_2| and \lstinline|uniq2|, but only if the user explicitly calls \lstinline|lem2|:
\begin{lstlisting}[caption=Equivalence Lemma for \lstinline|solution_2| and \lstinline|uniq2|]
lemma equivalence_solution2_uniq2(lst: List<int>)
  ensures (solution_2(lst) == uniq2(lst))
  decreases(length(lst))
{
  match lst
    case Nil => {}
    case Cons(hd, tl) =>
      lem2(hd, tl);
      equivalenceDrop(hd, tl);
}
\end{lstlisting}
%

\subsection{Equivalence of Imperative Methods}\label{sec:imperative-equivalence}

\footnote{Fix this missing lst:isSortedB reference}

Recursive functions can be converted into their iterative versions. For example, here is a recursive implementation of a function \lstinline|isSortedB| to determine whether a sequence of integers is sorted in non-decreasing order:
\begin{lstlisting}[caption=\lstinline|isSortedB| Function, label=lst:isSortedB]
function isSortedB(l: seq<int>): bool
  decreases |l|
{
  if |l| == 0 then true
  else if |l| == 1 then isSortedB([])
  else l[0] <= l[1] && isSortedB(l[1..])
}
\end{lstlisting}
%
This implementation can be converted into the following imperative method. It is not an exact conversion from the recursive version, to make better use of imperative features. The postcondition states that the method will output a boolean equivalent to whether each element in the sequence is less than or equal to the next element. If this can be proven, then the sequence must be sorted in non-decreasing order. The variable \lstinline|index| together with the invariant in the \lstinline|while| loop are used to keep track of which sequence elements have been checked so far.
\begin{lstlisting}[caption=\lstinline|isSortedBI| Method]
method isSortedBI(l: seq<int>) returns (res: bool)
  ensures res <==> (forall i :: 0 <= i < |l| - 1 ==> l[i] <= l[i+1])
{
  var index := 0;

  while index < |l| - 1
    invariant forall i :: 0 <= i < index ==> (i < |l| - 1 ==> l[i] <= l[i+1])
  {
    if l[index] > l[index + 1] {
        return false;
    }
    index := index + 1;
  }

  return true;
}
\end{lstlisting}
%
Proving the equivalence of imperative methods cannot be achieved directly through lemmas in Dafny; only functions can be called within lemmas as, by definition, Dafny functions are pure while Dafny methods are allowed to have side effects. Therefore, to prove equivalence, we must prove postconditions for both methods, and then state that these postconditions are equivalent. In turn, as in the above example, the user will often need to provide loop invariants and variants to help Dafny prove these postconditions.

These specifications can be very difficult to write correctly. Additionally, just proving that the methods' postconditions are equivalent is not sufficient to prove the equivalence of the methods themselves, as the postconditions may not fully capture the behaviour of the methods. 

\subsubsection{Proving Equivalence of Imperative Methods using Product Programs}

As described by \textcite{lahiriDifferentialAssertionChecking2013}, one solution to these difficulties is to generate product programs. A product program is a method that contains the bodies of the two methods being compared, in such a way that relationships between the two can be verified.

For example, we can use this process to prove the equivalence of the imperative versions of \lstinline|isSortedR| and \lstinline|isSortedB| (the function version of \lstinline|isSortedR| can be found in \Cref{lst:isSortedR}). The imperative version of \lstinline|isSortedR| is as follows, requiring a helper function \lstinline|loopI|. As above, the methods are not an exact conversion from their recursive counterparts. Additionally, the methods operate on sequences instead of lists, so I can use Dafny's built-in indexing and slicing features: 
\begin{lstlisting}[caption=\lstinline|isSortedRI| Method with \lstinline|loopI| Helper Method]
method loopI(p: int, l: seq<int>) returns (res: bool)
  ensures res <==> ((forall i :: 0 <= i < |l| - 1 ==> l[i] <= l[i+1]) && (|l| > 0 ==> p <= l[0]))
{
  var currentP := p;
  var next := 0;

  while next < |l|
    invariant next <= |l|
    invariant currentP == if next == 0 then p else l[next - 1]
    invariant next > 0 ==> p <= l[0]
    invariant forall i :: 0 <= i < next - 1 ==> l[i] <= l[i+1]
  {
    if currentP > l[next] {
        return false;
    }
    currentP := l[next];
    next := next + 1;
  }

  return true;
}

method isSortedRI(l: seq<int>) returns (res: bool)
  ensures res <==> (forall i :: 0 <= i < |l| - 1 ==> l[i] <= l[i+1])
{
  if |l| == 0 {
    return true;
  } else {
    var result := loopI(l[0], l[1..]);
    return result;
  }
}
\end{lstlisting}
%
These methods can be used to construct a product program as follows, using the definition of \lstinline|loopI| given above:
\begin{lstlisting}[caption=Product Program for \lstinline|isSortedBI| and \lstinline|isSortedRI| using \lstinline|loopI|]
method isSortedBI_RI_Equivalence(lb: seq<int>, lr: seq<int>) returns (rres: bool, bres: bool)
  requires lb == lr
  ensures rres == bres
{
  // Defining output variables
  var rout: bool;
  var bout: bool;

  // Computing isSortedRI
  if |lr| == 0 {
    rout := true;
  } else {
    var result := loopI(lr[0], lr[1..]);
    rout := result;
  }

  // Computing isSortedBI
  var index := 0;

  while index < |lb| - 1
    invariant forall i :: 0 <= i < index ==> (i < |lb| - 1 ==> lb[i] <= lb[i+1])
  {
    if lb[index] > lb[index + 1] {
        return rout, false;
    }

    index := index + 1;
  }

  bout := true;

  // Returning the results
  return rout, bout;
}
\end{lstlisting}
%

\footnote{Can you autogen these product programs? maybe without the invariants?}

The specification of this product program ensures that, since Dafny is able to verify this method, \lstinline|isSortedBI| and \lstinline|isSortedRI| must be equivalent. Therefore, this process removes the need for the user to provide postconditions for the individual methods, that capture their behaviour fully. In this case, a postcondition for the \lstinline|loopI| helper function is still required, but the body of this method can also be copied into the product program, removing this requirement as well. However, the user must still provide the loop invariants and variants for the individual methods.

\subsubsection{Equivalence of Methods and Functions}

A similar problem to that discussed above is proving the equivalence of a method and a function. In this case, the user does not need to write a postcondition for the method, since the function can be called directly within the specification. The user still needs to provide loop invariants and variants for the method, however, these do not need to be as complicated as those detailed above, as again, they can be defined in terms of the function. For example, the function \lstinline|isSortedB| defined in \Cref{lst:isSortedB} can be proven to be equivalent with the following imperative method:
\begin{lstlisting}[caption=\lstinline|isSortedBI| Method with simplified specifications]
method isSortedBI(l: seq<int>) returns (res: bool)
  ensures res == isSortedB(l)
{
  var index := 0;

  while index < |l| - 1
    invariant index <= |l|
    invariant isSortedB(l[index..]) == isSortedB(l)
  {
    if l[index] > l[index + 1] {
        return false;
    }
    index := index + 1;
  }

  return true;
}
\end{lstlisting}
%
Thus, the necessary specifications can be greatly simplified if they are defined in terms of the function the method is being compared to.

\subsection{Equivalence of Higher-Order Functions}\label{sec:higher-order}

Dafny is able to prove the equivalence of some higher-order functions, but not all. For example, here are two equivalent functions \lstinline|andThen1| and \lstinline|andThen2|:
\begin{lstlisting}[caption=\lstinline|andThen1| and \lstinline|andThen2| Functions]
function andThen1<A, B, C>(f: A -> B, g: B -> C): A -> C { a => g(f(a)) }
function andThen2<A, B, C>(ff: A -> B, gg: B -> C): A -> C { aa => gg(ff(aa)) }
\end{lstlisting}
%
Dafny is able to prove their equivalence without any help:
\begin{lstlisting}[caption=Equivalence Lemma for \lstinline|andThen1| and \lstinline|andThen2|]
lemma equivalenceAndThen<A, B, C>(f: A -> B, g: B -> C)
  ensures andThen1(f, g) == andThen2(f, g)
{}
\end{lstlisting}
%
Here are two more such functions \lstinline|flip1| and \lstinline|flip2|:
\begin{lstlisting}[caption=\lstinline|flip1| and \lstinline|flip2| Functions]
function flip1<A, B, C>(f: (A, B) -> C): (B, A) -> C {
  (b, a) => f(a, b)
}

function flip2<A, B, C>(f: (A, B) -> C): (B, A) -> C {
  (b, a) => var res := f(a, b); res
}
\end{lstlisting}
%
Dafny is able to prove that the outputted functions give the same output for the same inputs:
\begin{lstlisting}[caption=Equivalence Lemma for \lstinline|flip1| and \lstinline|flip2| with specified arguments]
lemma equivalenceFlip<A, B, C>(f: (A, B) -> C)
  ensures forall b: B, a: A :: flip1(f)(b, a) == flip2(f)(b, a)
{}
\end{lstlisting}
%
However, it is unable to prove that the outputted functions themselves are equivalent:
\begin{lstlisting}[caption=Invalid Equivalence Lemma for \lstinline|flip1| and \lstinline|flip2|]
// Dafny cannot prove this lemma
lemma equivalenceFlip<A, B, C>(f: (A, B) -> C)
  ensures flip1(f) == flip2(f)
{}
\end{lstlisting}
%
Therefore, Dafny is able to prove the equivalence of some higher-order functions. However, it struggles with others, despite being able to prove that the outputted functions give the same output for the same inputs.
